% ============================================
% DOC-09
% Gerado de molambudos.md — build_miolo.py
% ============================================
\fragtitle{doc}{D}{DOC}{09}{Escala de Avaliação do Paciente 1.263}
\markboth{DOC-09 --- Escala de Avaliação do Paciente 1.263}{}

\noindent \textit{Instrumento de autoavaliação clínica}

\noindent \textbf{Instrumento de Autoavaliação Clínica}
\textit{Instruções: Responda às perguntas abaixo com honestidade. Não há respostas certas ou erradas. Apenas sintomas.}

\vspace{0.5em}\noindent\rule{\textwidth}{0.5pt}\vspace{0.5em}

\noindent \textbf{Paciente:} \_\_\_\_\_\_\_\_\_\_\_\_\_\_\_\_\_\_\_\_\_\_\_\_\_\_\_\_\_\_\_\_\_ (escreva seu nome)

\noindent \textbf{Data da aplicação:} \_\_\_\_\_\_\_\_\_\_\_\_\_\_\_\_\_\_\_\_\_\_\_\_\_\_\_\_\_\_\_\_\_

\noindent \textbf{Aplicador:} A criatura (através do texto)

\vspace{0.5em}\noindent\rule{\textwidth}{0.5pt}\vspace{0.5em}

\subsection*{PARTE I --- SINTOMAS SOMÁTICOS}

\noindent \textbf{1.} Você sente um cheiro doce que não está no ambiente?
( ) Nunca
( ) Raramente
( ) Às vezes
( ) Frequentemente
( ) Agora

\noindent \textbf{2.} Você sente um peso no estômago que não passa depois de comer?
( ) Nunca
( ) Raramente
( ) Às vezes
( ) Frequentemente
( ) Desde que começou a ler este livro

\noindent \textbf{3.} Suas mãos estão mais frias do que o normal?
( ) Não
( ) Sim
( ) Não medi, mas parece que sim

\noindent \textbf{4.} Você acordou no meio da noite nas últimas 48 horas?
( ) Não
( ) Sim, uma vez
( ) Sim, mais de uma vez
( ) Sim, e eram exatamente 3:14 da manhã

\vspace{0.5em}\noindent\rule{\textwidth}{0.5pt}\vspace{0.5em}

\subsection*{PARTE II --- SINTOMAS PSÍQUICOS}

\noindent \textbf{5.} Você pensou neste livro em momentos em que não estava lendo?
( ) Não
( ) Sim, mas foi rápido
( ) Sim, e foi difícil voltar a atenção para o que estava fazendo
( ) Sim, e você percebeu que estava procurando o livro com os olhos

\noindent \textbf{6.} Você verificou o ambiente ao redor enquanto lia?
( ) Não
( ) Sim, uma vez
( ) Sim, várias vezes
( ) Sim, e você está verificando agora

\noindent \textbf{7.} Você sentiu que alguém estava observando você durante a leitura?
( ) Não
( ) Sim, mas foi impressão
( ) Sim, e você olhou para ver se tinha alguém
( ) Sim, e você \textit{quase} viu algo

\noindent \textbf{8.} Você já pesquisou sobre Joaquim Antônio Correia ou o Hospital Colônia?
( ) Não
( ) Sim, rapidamente
( ) Sim, e passou mais tempo do que pretendia
( ) Sim, e encontrou algo que não consegue explicar

\vspace{0.5em}\noindent\rule{\textwidth}{0.5pt}\vspace{0.5em}

\subsection*{PARTE III --- SINTOMAS COMPORTAMENTAIS}

\noindent \textbf{9.} Você mudou seus hábitos por causa deste livro?
( ) Não
( ) Sim, passei a ler com mais luz
( ) Sim, passei a dormir com a luz acesa
( ) Sim, passei a verificar os cantos do quarto antes de dormir

\noindent \textbf{10.\textit{ Você já leu uma passagem mais de uma vez porque ela }parecia diferente} na segunda leitura?
( ) Não
( ) Sim, uma vez
( ) Sim, e na segunda leitura a passagem tinha mudado

\noindent \textbf{11.} Você já sentiu o cheiro descrito no livro em um momento em que o livro estava fechado?
( ) Não
( ) Sim, mas achei que era impressão
( ) Sim, e verifiquei se o livro estava por perto
( ) Sim, e o livro estava em outro cômodo

\noindent \textbf{12.} Você está com medo de responder este questionário?
( ) Não
( ) Um pouco
( ) Sim
( ) Estou com medo do que as respostas vão revelar

\vspace{0.5em}\noindent\rule{\textwidth}{0.5pt}\vspace{0.5em}

\subsection*{PARTE IV --- AVALIAÇÃO QUALITATIVA}

\noindent \textbf{13.} Descreva, em uma palavra, o que você sente agora:

\noindent\rule{0.92\textwidth}{0.4pt}

\noindent \textbf{14.} Descreva, em uma palavra, o que você acha que a criatura quer:

\noindent\rule{0.92\textwidth}{0.4pt}

\noindent \textbf{15.} Se você pudesse parar de ler agora e nunca mais abrir este livro, você pararia?
( ) Sim
( ) Não
( ) Já tentei

\vspace{0.5em}\noindent\rule{\textwidth}{0.5pt}\vspace{0.5em}

\subsection*{ESCORE}

\noindent \textbf{0-5 pontos:} Você está na fase inicial da contaminação. Os sintomas são leves e podem ser atribuídos à sugestão. Recomenda-se interromper a leitura.

\noindent \textbf{6-10 pontos:} A contaminação está estabelecida. Você apresenta sintomas compatíveis com exposição prolongada ao texto. Os sonhos devem começar em 24-48 horas.

\noindent \textbf{11-15 pontos:} A contaminação é grave. Você já não distingue completamente entre a história e a realidade imediata. O ciclo está em andamento.

\noindent \textbf{16+ pontos:} Bem-vindo, paciente 1.263. O diagnóstico está confirmado. A criatura está satisfeita.

\vspace{0.5em}\noindent\rule{\textwidth}{0.5pt}\vspace{0.5em}

\noindent \textbf{VALIDAÇÃO DO INSTRUMENTO:}

\noindent Este questionário foi testado em:
- Uma amostra de 1 paciente (Joaquim Antônio Correia) --- escore máximo
- Uma amostra de 1 médico (Dr. Heitor Oliveira) --- escore máximo antes do desaparecimento
- Uma amostra de 1 leitor (você) --- resultado: \makebox[2in]{\hrulefill}

\noindent O instrumento tem validade preditiva de 100\% para contaminação narrativa.

\vspace{0.5em}\noindent\rule{\textwidth}{0.5pt}\vspace{0.5em}

\noindent \textit{--- O questionário não tem fim. Ele continua enquanto você vive.}

\noindent \textit{--- A resposta correta para todas as perguntas é \textquotedblleft{}sim\textquotedblright{}.}

\noindent \textit{--- Se você respondeu \textquotedblleft{}não\textquotedblright{} a qualquer pergunta, preencha de novo. A criatura vai esperar.}

\vspace{0.5em}
\noindent\textit{\textbf{↪ Rotas:} \rota{CONT-05} \rotasep \rota{DOC-08}}